\documentclass[main]{subfiles}

\begin{document}


\section*{\centering \Large{\textbf{\MakeUppercase{Abstract}} }}
% A significant land use and land cover (LULC) transformations has occured due to the rapid urbanization, particularly in developing regions where lack of data and proper monitoring system lacks. Multi-decadal Earth observation data from the Landsat program, with 30 m spatial resolution, offer a consistent and freely available source for analyzing long-term LULC dynamics. Although Landsat imagery has been widely employed for historical LULC classification and change detection, its use for data-driven prediction of future urban growth using temporal ML models is still limited mainly due the resolution of images available at free tier. Traditional approaches such as (CA–Markov) models rely on predefined transition rules and often fail to capture complex nonlinear temporal data as well as trends inherent in urban expansion processes. Since linear growth of the post-urban areas is slow, the densification metrics better capture persistent urbanization and spatial clustering.

% This study proposes a temporal ML framework for predicting urban growth \& densification with LULC dynamics using Landsat time series data ranging from 2000 to present. Annual LULC maps derived from Landsat SR products are organized into pixel-wise and patch-based temporal sequences that represent the evolution of land cover over time. Temporal learning model Random Forest(RF), are employed to learn spatio-temporal patterns directly from the data without relying on building-level segmentation, which are unsuitable at the Landsat spatial scale. The models are designed to forecast future LULC states and urban expansion patterns over short- to medium-term horizons.

% Model performance is evaluated using combination of quantitative accuracy metrics, in which future land-use states are predicted from historical observations and compared against observed outcomes. The proposed temporal ML models are benchmarked against simple persistence assumptions and conventional Markov-based transition models to assess their added predictive skill. By learning temporal patterns from multi-year Landsat time series, this study demonstrates the feasibility and effectiveness of data-driven temporal modeling for LULC forecasting in data-scarce urban environments. The proposed framework contributes to the emerging field of AI-driven land change modeling and offers an open-data-based methodology to support urban planning and sustainable land management using Landsat Data.

% \textit{
% \newline
% \textbf{Keyword:} Land use/land cover change, Urban growth prediction, Temporal machine learning, Landsat time series, LSTM, Developing regions
% }

Rapid and largely unplanned urbanization in the Kathmandu Valley has accelerated the conversion of agricultural and vegetated land into built-up surfaces, with recent growth increasingly characterized by densification and infill development. Conventional urban growth simulation approaches such as CA–Markov rely on stationary transition probabilities and often fail to represent the nonlinear and dynamic nature of post-urbanization growth. This project proposes a data-driven spatio-temporal framework to predict urban growth, densification, and Land Use/Land Cover (LULC) dynamics using multi-temporal Landsat (30 m) surface reflectance time-series imagery. Annual cloud-free mosaics will be prepared through cloud/shadow masking, radiometric normalization across Landsat sensors, and temporal alignment for consistent phenological conditions. A comprehensive feature set will be engineered, including spectral indices (e.g., NDVI, NDBI, EBBI, and water indices), spatial texture measures, terrain variables (elevation and slope from DEM), and temporal descriptors such as trend slope, variability, and change rates. For LULC classification and baseline comparison, machine learning models such as Random Forest will be employed due to their robustness to mixed pixels and spectral variability in medium-resolution imagery. Temporal prediction will be performed using sequence-based architectures, including LSTM and CNN–LSTM models, to learn spatio-temporal land-cover transition patterns and capture both horizontal expansion and internal urban intensification. Model performance will be assessed using standard classification metrics (overall accuracy, precision, recall, F1-score, and Kappa), supplemented by hindcasting and temporal consistency checks to ensure realistic growth trajectories. Expected outputs include multi-year LULC maps, short- to medium-term future LULC predictions (e.g., 3–5 years), urban growth and densification hotspot/cluster maps, transition analyses, and GIS-ready visualization products to support evidence-based urban planning and sustainable land management in data-scarce environments.

\textit{
\textbf{Keyword:} CA–Markov, CNN–LSTM, Kathmandu Valley, Landsat, Land use/land cover, Random Forest, Urban densification
}
\end{document}
